\documentclass[11pt]{article}

\usepackage[margin=1in]{geometry}
\usepackage{amsmath,amssymb,amsthm,booktabs,mathtools}
\usepackage[numbers]{natbib}
\usepackage[colorlinks=true,allcolors=blue]{hyperref}
\usepackage[nameinlink,noabbrev]{cleveref}

\newtheorem{proposition}{Proposition}
\newtheorem{remark}{Remark}

\newcommand{\StageThreeSectionNote}{}
\newcommand{\SupplementTheoryOnlyNote}{}
\newcommand{\SupplementMixedEvidenceNote}{}

\title{Coordination Persistence Flagship Paper: Working Draft}
\author{}
\date{\today}

\begin{document}

\maketitle

\begin{abstract}
This standalone repo carries forward only the locked Stage 3 theory and methods
surfaces from the parent project. Fresh data collection, fresh feature
construction, and all new empirical results will be produced inside this repo.
\end{abstract}

\section{Scope note}

This working draft intentionally excludes legacy results, discussion, and
conclusion text. The purpose of the current paper wrapper is to preserve the
locked theoretical and methodological core while the new flagship evidence base
is built from fresh data.

\section{Framework and Estimands}
\label{sec:model}
\StageThreeSectionNote

This paper studies coordination persistence as a risk-set-defined continuation
problem on non-root comments in Moltbook. Each non-root parent comment creates a
follow-up opportunity: another comment either arrives as a direct reply within
an operational horizon or it does not. The central OR/MS question is therefore
not whether activity exists somewhere on the platform, but which margin limits
within-horizon continuation on a common opportunity set---reply incidence or
conditional reply speed. The framework below defines that control panel, links
it to thread-level coordination outcomes, and states the exact algebraic
propositions that discipline estimation and intervention-priority logic.

\subsection{Scientific object and scope}
\label{sec:model:scope}

Let \(m\) index a non-root comment that can itself receive a direct reply.
Write \(\tau_m\) for its creation time and \(T_m\) for the lag from \(\tau_m\)
to the first direct reply, if one occurs. Let \(C_m\) denote observed follow-up
from \(\tau_m\) until administrative end, archive interruption, or other
coverage loss, and define
\[
\delta_m := \mathbf{1}\{T_m \le C_m\}.
\]
All primary incidence and timing estimands are defined on these non-root parent
comments; direct replies to root posts are outside the primary persistence
object. This scope isolates continuation beyond the first layer and avoids
conflating root attraction with sustained comment-to-comment coordination.

The paper is observational throughout. It measures continuation and its
structural consequences in archived interaction traces, but it does not identify
causal intervention effects, latent exposure, or pure autonomous intent. These
boundaries are important because the same observed reply pattern can be
generated by multiple latent mechanisms, including resurfacing, visibility,
moderation, memory loss, or scheduled agent availability
\citep{hawkes1971spectra,whitt2004efficiency,reed2012hazard,jouini2010online}.

\subsection{Observed thread state and downstream coordination outcomes}
\label{sec:model:branching}

Let \(j \in \mathcal{J}\) index a thread with root post \(n=0\) and comments
\(n=1,\ldots,N_j\). Event \(n\) is \((t_{jn},a_{jn},p_{jn})\), where
\(t_{jn}\) is the UTC timestamp, \(a_{jn}\) is the observed author identifier,
and \(p_{jn}\in\{0,\ldots,n-1\}\) is the parent index. Node depth is defined
recursively by
\begin{equation}
\label{eq:depth}
d_{j0}=0,
\qquad
d_{jn}=d_{j,p_{jn}}+1 \quad (n\ge 1),
\end{equation}
and maximum thread depth is \(D_j:=\max_{0\le n\le N_j} d_{jn}\).

For node \((j,n)\), let
\[
c_{jn}:=\#\{r>n : p_{jr}=n\}
\]
be its direct-child count. The branching-by-depth profile is
\begin{equation}
\label{eq:branching-by-depth}
\bar c_k := \mathbb{E}\!\left[c_{jn}\mid d_{jn}=k\right], \qquad k\ge 0.
\end{equation}
This profile distinguishes root-heavy star patterns from deeper continuation.

Define the directed reply-edge set
\[
\begin{aligned}
E_j
&:=\{(u,v): \exists n\ge 1 \text{ with } a_{jn}=u,\\
&\hspace{2.7em} a_{j,p_{jn}}=v,\ u\neq v\},
\end{aligned}
\]
and let
\[
\begin{aligned}
\Delta_j
&:=\{\{u,v\} : (u,v)\in E_j\\
&\hspace{2.7em} \text{or } (v,u)\in E_j\}
\end{aligned}
\]
be the unordered dyads with at least one directional reply. Thread-level
reciprocity is
\begin{equation}
\label{eq:reciprocity-rate}
R_j
:=\frac{1}{|\Delta_j|}\sum_{\{u,v\}\in \Delta_j}
\mathbf{1}\{(u,v)\in E_j \text{ and } (v,u)\in E_j\},
\end{equation}
with \(R_j\) left undefined when \(|\Delta_j|=0\).

Thread-level re-entry is
\begin{equation}
\label{eq:reentry-rate}
RE_j
:=\frac{\#\{n\ge 1 : a_{jn}\in\{a_{j1},\ldots,a_{j,n-1}\}\}}{N_j},
\end{equation}
again defined on the non-root comment sequence. Together,
\[
G_j := \left(D_j,\{\bar c_k\}_{k\ge 0},R_j,RE_j\right)
\]
summarizes the downstream coordination structure that accumulates from
micro-level continuation successes and failures
\citep{harris1963theory,gomez2013structure,aragon2017thread,meital2024branch}.

\subsection{The horizon-throughput control panel}
\label{sec:model:estimands}

For horizon \(h>0\), define the horizon risk-set indicator
\[
R_h(m):=\mathbf{1}\{C_m\ge h \text{ or } T_m\le h\}.
\]
The paper's primary object is the common-risk-set control panel
\[
\mathcal{C}_h := (\pi_h,\phi_h,q_h), \qquad h\in\{5\text{ min},1\text{ h}\},
\]
where
\begin{align}
\pi_h &:= \Pr(\delta_m=1 \mid R_h(m)=1), \\
\phi_h &:= \Pr(T_m\le h \mid \delta_m=1,\ R_h(m)=1), \\
q_h &:= \Pr(T_m\le h \mid R_h(m)=1).
\end{align}
Earlier drafts reported \(q_h\) under the notation \(p_h\); throughout the
final architecture, \(q_h\) is the preferred notation because it makes the
incidence--timing decomposition explicit.

The secondary descriptive incidence margin is
\[
p_{\mathrm{obs}}:=\Pr(\delta_m=1),
\]
the coverage-conditional ever-reply share. It is reported for descriptive
orientation only and is not the flagship decision object because it is sensitive
to heterogeneous follow-up.

Conditional timing is summarized nonparametrically through the distribution of
\(T_m\mid (\delta_m=1)\), especially
\((t_{50},t_{90},t_{95})\) and short-window masses such as
\(\Pr(T_m\le 30\text{ s}\mid \delta_m=1)\) and
\(\Pr(T_m\le 5\text{ min}\mid \delta_m=1)\).

\begin{proposition}[Common-risk-set factorization]
\label{prop:risk-set-factorization}
For any horizon \(h\) with \(\Pr(R_h(m)=1)>0\),
\begin{equation}
\label{eq:horizon-throughput-factorization}
q_h = \pi_h \phi_h.
\end{equation}
\end{proposition}

\begin{proof}
See Supplementary Material, Proposition 1.
\end{proof}

\subsection{Minimal mechanism scaffold and identification boundary}
\label{sec:model:mechanism}

The paper retains only a minimal observational mechanism scaffold. For parent
comment \(m\) and age \(u>0\), write
\begin{equation}
\label{eq:minimal-mechanism}
\lambda_m(u)=E_m(u)\,S_m(u),
\end{equation}
where \(E_m(u)\) is effective exposure or availability and \(S_m(u)\) is a
weakly decreasing staleness kernel. Heartbeat routines, resurfacing, ranking,
and platform visibility all load into \(E_m(u)\); memory decay, context loss,
and reply urgency load into \(S_m(u)\). This factorization is deliberately
agnostic: the current archives do not separately identify exposure from
staleness, so \Cref{eq:minimal-mechanism} is used as disciplined mechanism
vocabulary rather than as a fully estimated structural model
\citep{hawkes1971spectra,crane2008robust,cox1972regression}.

A parametric exponential kernel \(S_m(u)=\exp(-\beta u)\) is therefore only a
secondary timing diagnostic. When such a diagnostic is reported, the quantity
\(\ln 2/\beta\) is interpreted as an exponential-equivalent reply-hazard
timescale, not as a thread ``half-life'' and not as a primary estimand.

\subsection{Design assumptions and primary hypotheses}
\label{sec:model:assumptions}

The framework rests on the following design assumptions.

\begin{description}
\item[A1. Archive integrity and non-root scope.]
After deterministic deduplication and UTC normalization, parent links and event
ordering are accurate enough for aggregate tree reconstruction, and the primary
persistence object is restricted to non-root comments.

\item[A2. Risk-set identification and gap discipline.]
For fixed horizon \(h\), the quantities \(q_h\), \(\pi_h\), and \(\phi_h\) are
identified from parents with at least \(h\) observed follow-up or an observed
reply before \(h\). Coverage gaps are treated as missing exposure windows rather
than filled by imputation.

\item[A3. Cross-archive comparability.]
A larger public Moltbook archive can be restricted to a harmonized minimal
schema---thread identifier, comment identifier, parent identifier, author
identifier, UTC timestamp, and community label---sufficient for separate-archive
replication of the supported core metrics.

\item[A4. Mechanism non-separability.]
Without impressions, ranking, or exposure logs, visibility/availability effects
and staleness/memory effects are not separately identified from archival traces.

\item[A5. Local comparative statics.]
The intervention logic concerns local or one-margin perturbations to the
observed margins \((\pi_h,\phi_h)\); the archival data do not identify the
causal intervention effects that would generate those perturbations.

\item[A6. Coordination bridge.]
Higher sustained non-root horizon throughput should weakly support deeper
continuation, more non-root branching, more reciprocity, and more re-entry, but
the exact branching law is not identified.

\item[A7. Coarse heterogeneity proxies.]
Deterministic submolt labels are adequate for prespecified coarse heterogeneity,
whereas account-status fields are descriptive only and do not support mechanism
attribution.

\item[A8. Observed-dynamics and periodicity boundary.]
The paper makes claims about observed platform interaction dynamics only.
Aggregate periodicity diagnostics can detect strong synchronization when it is
present, but weak concentration does not rule out individual routines under
dephasing or jitter.
\end{description}

The empirical design tests the following hypotheses.

\begin{description}
\item[H1 (core).]
After gap-aware and horizon-standardized estimation, Moltbook exhibits a stable
low-incidence / fast-conditional regime on non-root reply opportunities at
\(h=5\) minutes and \(h=1\) hour. The same ordering is expected to survive
separate-archive replication on a larger public Moltbook archive once the
required minimal schema is verified.

\item[H2 (core).]
At minute-to-hour horizons, Moltbook is incidence-constrained rather than
timing-constrained: \(\phi_h\) materially exceeds \(\pi_h\), so equal-size or
moderate-cost-ratio perturbations imply larger local throughput gains from
incidence lift than from timing lift.

\item[H3 (core).]
Windows or strata with lower \(q_h\) also exhibit shallower depth tails, lower
non-root branching, lower reciprocity, and lower re-entry, consistent with
coordination-persistence limits.

\item[H4 (secondary).]
Prespecified coarse topic strata differ materially in \(q_h\) and in upper-tail
conditional timing, and the main direction of those contrasts is robust to
deterministic keyword-remapping sensitivity.

\item[H5 (diagnostic only).]
Aggregate four-hour periodicity is weak and non-load-bearing: departures from
exact phase uniformity, if detected, are too small to serve as the main
explanation of the persistence bottleneck.
\end{description}

\subsection{Operational decision rule}
\label{sec:model:or-diagnostic}

Because \Cref{eq:horizon-throughput-factorization} is exact, local one-margin
perturbations yield exact first-step gains. Let \(\Delta q_h^{(\pi)}\) denote
the change in \(q_h\) when only \(\pi_h\) changes by \(\Delta \pi_h\), holding
\(\phi_h\) fixed, and define \(\Delta q_h^{(\phi)}\) analogously when only
\(\phi_h\) changes.

\begin{proposition}[One-margin gains and cost-normalized priority]
\label{prop:local-gain-priority}
If only incidence changes, then
\[
\Delta q_h^{(\pi)}=\phi_h\,\Delta \pi_h.
\]
If only conditional timing changes, then
\[
\Delta q_h^{(\phi)}=\pi_h\,\Delta \phi_h.
\]
For positive intervention costs \(c_\pi\) and \(c_\phi\), the corresponding
cost-normalized local gains are
\begin{equation}
\label{eq:horizon-throughput-delta}
I_\pi:=\frac{\phi_h\,\Delta \pi_h}{c_\pi},
\qquad
I_\phi:=\frac{\pi_h\,\Delta \phi_h}{c_\phi},
\end{equation}
so local priority favors incidence lift if and only if \(I_\pi>I_\phi\).
\end{proposition}

\begin{proof}
See Supplementary Material, Proposition 2.
\end{proof}

\begin{remark}[Auxiliary coordination bridge]
The paper tests the link from \(q_h\) to downstream thread structure through
joint readouts of \((q_h,D_j,\bar c_k,R_j,RE_j)\). No identified branching
ratio or exact reproduction law is claimed. Any approximation such as
\(\Pr(D_j\ge K)\approx q_h^{K-1}\) is therefore heuristic and explicitly
non-load-bearing.
\end{remark}

\section{Identification, Estimation, and Empirical Procedures}
\label{sec:methods}
\StageThreeSectionNote

This study is an observational measurement paper with disciplined theory rather
than a causal intervention study. The identified objects are the common-risk-set
margins \((q_h,\pi_h,\phi_h)\), the conditional timing summaries among observed
replies, and the thread-geometry outcomes defined in \Cref{sec:model}. Exposure
and staleness are not separately identified, and the policy indices in
\Cref{sec:model:or-diagnostic} are plug-in decision calculations on observed
margins rather than estimates of intervention effects.

\subsection{Archive discipline and current evidence boundary}
\label{sec:methods:design}

The canonical parent-linked archive is the public SimulaMet Moltbook snapshot
\citep{simulamet2026observatoryarchive}. The analysis begins with deterministic
deduplication, UTC normalization, and parent-link validation. Each non-root
comment becomes one candidate-parent unit, and each root post becomes one thread
for the geometry readouts. The present manuscript therefore reports only
SimulaMet-backed parent-linked evidence.

The study design also includes a separate larger-archive replication on a public
Moltbook archive (the approved MoltNet replication surface) under a harmonized
minimal schema. That replication is not yet implemented in the current evidence
base, so no pooled cross-archive estimator or cross-archive average is reported
here. Until minimal-schema equivalence is audited, archive disagreement is
treated as a replication object rather than pooled away.

\subsection{Primary risk-set estimators}
\label{sec:methods:def-est}
\label{sec:methods:two-part}

Let \(\mathcal{M}\) denote the realized cohort of non-root parent comments. For
parent \(m\in\mathcal{M}\) and horizon \(h\), define the realized within-horizon
completion indicator
\[
Y_h(m):=\mathbf{1}\{\delta_m=1,\ T_m\le h\}
\]
and the horizon risk-set indicator
\[
R_h(m):=\mathbf{1}\{C_m\ge h \text{ or } T_m\le h\}.
\]
The primary horizons are \(h\in\{5\text{ min},1\text{ h}\}\), with
\(30\) seconds used only in gap and censoring diagnostics.

The sample analogues of the control panel in \Cref{sec:model:estimands} are
\begin{align}
\widehat q_h
&=
\frac{\sum_{m\in\mathcal{M}} Y_h(m)}
     {\sum_{m\in\mathcal{M}} R_h(m)},
\\
\widehat \pi_h
&=
\frac{\sum_{m\in\mathcal{M}} \mathbf{1}\{\delta_m=1\}R_h(m)}
     {\sum_{m\in\mathcal{M}} R_h(m)},
\\
\widehat \phi_h
&=
\frac{\sum_{m\in\mathcal{M}} Y_h(m)}
     {\sum_{m\in\mathcal{M}} \mathbf{1}\{\delta_m=1\}R_h(m)}.
\end{align}
By construction, the sample identity
\[
\widehat q_h=\widehat \pi_h\,\widehat \phi_h
\]
holds exactly whenever the three quantities are computed on the same risk set.
Earlier draft tables denoted \(\widehat q_h\) by \(\widehat p_h\); the present
section uses \(\widehat q_h\) to reserve \(\widehat \pi_h\) and
\(\widehat \phi_h\) for its two interpretable margins.

The secondary descriptive incidence measure is
\[
\widehat p_{\mathrm{obs}}
=
\frac{1}{|\mathcal{M}|}\sum_{m\in\mathcal{M}} \mathbf{1}\{\delta_m=1\},
\]
the in-window ever-replied share. Because it is sensitive to heterogeneous
follow-up, \(\widehat p_{\mathrm{obs}}\) is treated as descriptive orientation
only and not as the paper's main decision object.

Among replied parents, conditional timing is summarized nonparametrically by the
empirical distribution of \(T_m\mid (\delta_m=1)\), especially
\((t_{50},t_{90},t_{95})\), plus short-window masses such as
\(\Pr(T_m\le 30\text{ s}\mid \delta_m=1)\) and
\(\Pr(T_m\le 5\text{ min}\mid \delta_m=1)\). Their unconditional counterparts
\(\Pr(\delta_m=1,\ T_m\le t)\) are reported on the same parent cohort.

\subsection{Gap handling, horizon standardization, and archive-conflict discipline}
\label{sec:methods:gap-discipline}

The canonical comment stream contains a 41.7-hour interruption, so gap handling
is part of the primary design rather than a late robustness appendix. The paper
therefore treats unobserved intervals as missing coverage rather than as true
silence and uses three built-in checks.

First, the core control-panel quantities are recomputed on contiguous windows,
including a short pre-gap window and the substantially longer post-gap window.
Second, parent comments whose potential follow-up windows overlap the gap are
excluded in prespecified overlap checks. Third, the headline completion
quantities are reported at explicit horizons using the risk-set estimators above,
so right-censoring near the observation boundary cannot masquerade as slow
conditional timing.

No replies are imputed across the gap. Likewise, if a later archive conflicts
with the canonical snapshot because of schema or coverage differences, that
conflict is handled through separate re-estimation rather than by pooled
cross-archive fusion. This is why gap handling, horizon standardization, and
archive-conflict discipline are treated as core identification design choices.

\subsection{Conversation geometry and downstream coordination readouts}
\label{sec:methods:geometry}

For each thread \(j\), depth is computed from \Cref{eq:depth}. The primary
geometry readouts are the empirical depth-tail probabilities
\(\Pr(D_j\ge k)\), maximum depth \(D_j\), and the branching-by-depth profile
\(\bar c_k\) from \Cref{eq:branching-by-depth}. We additionally summarize
reciprocity \(R_j\) and re-entry \(RE_j\) using
\Cref{eq:reciprocity-rate,eq:reentry-rate}. These outcomes are not treated as
independent mechanisms; instead, they provide the downstream coordination bridge
for testing whether lower \(q_h\) coincides with shallower continuation and
weaker repeated participation.

Where a compact depth-tail summary is useful, the paper reports an effective
log-tail slope \(\hat s_{\mathrm{depth}}\) fitted to
\(\log \Pr(D_j\ge k)\) for \(k\ge 2\). That slope is descriptive. It is not
claimed to be an identified reproduction parameter.

\subsection{Auxiliary stratified and parametric diagnostics}
\label{sec:methods:auxiliary}

Submolt heterogeneity is evaluated on the same risk-set estimands
\((\widehat q_h,\widehat \pi_h,\widehat \phi_h)\) using prespecified
deterministic topic strata. Topic heterogeneity is secondary; it remains useful
only insofar as it sharpens which persistence margin is more constrained.
Account-status contrasts are retained as descriptive heterogeneity only and do
not carry mechanism or causal interpretation.

For auxiliary descriptive incidence contrasts across coarse strata, we fit a
complementary log-log incidence regression:
\begin{equation}
\label{eq:methods-incidence-cloglog}
\log\!\left[-\log\!\left(1-\Pr(\delta_m=1\mid x_m)\right)\right]
=
x_m^{\top}\eta.
\end{equation}
Here \(x_m\) contains prespecified group indicators such as submolt class and
account-status group. This regression is calibration-oriented: it summarizes
direction and relative incidence differences, but it is not the paper's primary
estimand.

Parametric timing fits are also auxiliary only. For diagnostic comparison with
the nonparametric timing summaries, we fit an exponential-kernel hazard model
and a Weibull alternative. The exponential-kernel log-likelihood is
\begin{equation}
\label{eq:exponential-ll}
\ell(\alpha,\beta)
=
\sum_m
\left[
\delta_m(\log \alpha-\beta s_m)
-
\frac{\alpha}{\beta}\left(1-e^{-\beta s_m}\right)
\right],
\end{equation}
where \(s_m:=\min(T_m,C_m)\), and the Weibull survival benchmark is
\begin{equation}
\label{eq:weibull-survival}
S(s)=\exp\!\left[-\left(\frac{s}{\lambda}\right)^{\gamma}\right].
\end{equation}
These fits are used only as misspecification diagnostics and as secondary
timescale summaries; they do not determine the main conclusions.

\begin{remark}[Kernel half-life diagnostic]
\label{rem:estimand}
The quantity \(\ln 2/\beta\) from the exponential timing diagnostic is an
exponential-equivalent reply-hazard timescale. It is not a median thread
lifetime, and it is not a substitute for the primary control-panel margins.
\end{remark}

The repository analysis pipeline also reports a boundary-exclusion sensitivity
for the half-life diagnostic that removes parents created near the observation
end. That sensitivity is supplementary only and does not alter the primary
risk-set estimands.

\subsection{Diagnostic periodicity layer}
\label{sec:methods:periodicity}

Heartbeat-style synchronization is assessed only diagnostically and only on the
longest contiguous segment of the comment stream. The diagnostic stack includes:
event-time modulo-4-hour concentration, the Rayleigh statistic, Monte Carlo
\(p\)-values for exact phase uniformity, power spectral density (PSD) checks at
the four-hour target frequency, first-order autoregressive (AR(1)) null
calibration for binned-count spectra, and bin-width sensitivity.

This layer is intentionally non-load-bearing. At large \(N\), exact uniformity
can be rejected even when concentration is far too small to imply practically
meaningful global synchronization. Conversely, weak aggregate concentration does
not rule out individual routines under dephasing or jitter. Periodicity is
therefore treated as mechanism context only.

\subsection{Separate-archive replication requirement}
\label{sec:methods:replication}

The study design extends beyond the canonical first-week archive. The same
supported core metrics are to be re-estimated on a larger public Moltbook
archive under a harmonized minimal schema containing thread identifier, comment
identifier, parent identifier, author identifier, UTC timestamp, and community
label. Replication is performed separately by archive. No pooled cross-archive
estimator is used unless schema equivalence and coverage comparability have been
audited first.

If the replication archive preserves reliable parent-comment linkage, the same
geometry outcomes \((D_j,\bar c_k,R_j,RE_j)\) can be transported as well. If
parent linkage is incomplete or unreliable, the replication narrows to the
supported two-part persistence quantities \((q_h,\pi_h,\phi_h)\) and the
conditional timing summaries. In the current manuscript evidence base, however,
only the canonical SimulaMet archive supports the reported parent-linked
estimates.

\subsection{Comparator boundary and uncertainty}
\label{sec:methods:comparison}
\label{sec:methods:uncertainty}

No human-platform comparator enters the paper's primary estimands,
identification logic, or policy index. Any residual comparator analysis kept
elsewhere in the manuscript should be read as descriptive context only and not
as part of the flagship architecture.

Uncertainty for the nonparametric control-panel and geometry summaries is
reported with thread-cluster bootstrap intervals using fixed resampling rules.
Auxiliary incidence regressions use cluster-aware covariance that accounts for
dependence within threads and repeated authors where those identifiers are
available. One-parent-per-thread sensitivity is used as a further dependence
check. These procedures quantify uncertainty around observational readouts; they
do not convert the study into a causal design.


\appendix
\noindent This supplementary material records the exact control-panel proofs,
the estimator definitions that bridge theory to observables, and the gap-aware
diagnostic tables retained in the approved manuscript architecture. It
deliberately omits comparator-centered appendices and theorem-forward material
from earlier drafts because those objects are not load-bearing in the final
design.

\section*{S1. Exact control-panel propositions}
\SupplementTheoryOnlyNote

The main text relies on only two exact propositions. They are restated and
proved here.

\begin{proposition}[Common-risk-set factorization]
Fix a horizon \(h>0\). Define
\[
\begin{aligned}
R_h(m)
&:=\mathbf{1}\!\left\{\substack{C_m\ge h\\ \text{or } T_m\le h}\right\},\\
\pi_h
&:=\mathbb{P}(\delta_m=1\mid R_h(m)=1),\\
\phi_h
&:=\mathbb{P}\!\left(T_m\le h \,\middle|\, \substack{\delta_m=1\\ R_h(m)=1}\right),\\
q_h
&:=\mathbb{P}(T_m\le h \mid R_h(m)=1).
\end{aligned}
\]
If \(\mathbb{P}(R_h(m)=1)>0\) and
\(\mathbb{P}(\delta_m=1\mid R_h(m)=1)>0\), then
\[
q_h=\pi_h\phi_h.
\]
\end{proposition}

\begin{proof}
Conditional on \(R_h(m)=1\), the event \(\{T_m\le h\}\) implies
\(\{\delta_m=1\}\): if the first direct reply arrives by \(h\), then a direct
reply is observed before censoring. Therefore,
\begin{align*}
\mathbb{P}(T_m\le h\mid R_h(m)=1)
&=
\mathbb{P}(T_m\le h,\ \delta_m=1 \\
&\hspace{2.5em}\mid R_h(m)=1).
\end{align*}
Applying the chain rule on the right-hand side gives
\begin{gather*}
\mathbb{P}(T_m\le h,\ \delta_m=1\mid R_h(m)=1)\\
=\mathbb{P}(\delta_m=1\mid R_h(m)=1)\times \\
\mathbb{P}\!\left(T_m\le h \,\middle|\, \substack{\delta_m=1\\ R_h(m)=1}\right).
\end{gather*}
which is exactly \(q_h=\pi_h\phi_h\).
\end{proof}

\begin{proposition}[One-margin gains and cost-normalized priority]
Let \(q_h=\pi_h\phi_h\). If only the incidence margin changes by
\(\Delta \pi_h\) while \(\phi_h\) is held fixed, then
\[
\Delta q_h^{(\pi)}=\phi_h\,\Delta \pi_h.
\]
If only the conditional-timing margin changes by \(\Delta \phi_h\) while
\(\pi_h\) is held fixed, then
\[
\Delta q_h^{(\phi)}=\pi_h\,\Delta \phi_h.
\]
For positive intervention costs \(c_\pi\) and \(c_\phi\), the corresponding
cost-normalized local gains are
\[
I_\pi:=\frac{\phi_h\,\Delta \pi_h}{c_\pi},
\qquad
I_\phi:=\frac{\pi_h\,\Delta \phi_h}{c_\phi},
\]
so local priority favors incidence lift if and only if \(I_\pi>I_\phi\).
\end{proposition}

\begin{proof}
If only incidence changes, then
\[
q_h+\Delta q_h^{(\pi)}
=
(\pi_h+\Delta \pi_h)\phi_h
=
\pi_h\phi_h+\phi_h\Delta \pi_h,
\]
hence \(\Delta q_h^{(\pi)}=\phi_h\Delta \pi_h\). If only conditional timing
changes, then
\[
q_h+\Delta q_h^{(\phi)}
=
\pi_h(\phi_h+\Delta \phi_h)
=
\pi_h\phi_h+\pi_h\Delta \phi_h,
\]
hence \(\Delta q_h^{(\phi)}=\pi_h\Delta \phi_h\). Dividing by positive costs
\(c_\pi\) and \(c_\phi\) yields the indices \(I_\pi\) and \(I_\phi\), so the
preferred local margin is incidence if and only if \(I_\pi>I_\phi\).
\end{proof}

\section*{S2. Observable bridge and primary estimators}
\SupplementMixedEvidenceNote

Let \(\mathcal{M}\) denote the realized cohort of non-root parent comments.
For \(m\in\mathcal{M}\) and horizon \(h\), define
\[
\begin{aligned}
Y_h(m)&:=\mathbf{1}\{\delta_m=1,\ T_m\le h\},\\
R_h(m)&:=\mathbf{1}\!\left\{\substack{C_m\ge h\\ \text{or } T_m\le h}\right\}.
\end{aligned}
\]
The sample control-panel estimators are
\begin{align*}
\widehat q_h
&=
\frac{\sum_{m\in\mathcal{M}} Y_h(m)}
     {\sum_{m\in\mathcal{M}} R_h(m)},
\\
\widehat \pi_h
&=
\frac{\sum_{m\in\mathcal{M}} \mathbf{1}\{\delta_m=1\}R_h(m)}
     {\sum_{m\in\mathcal{M}} R_h(m)},
\\
\widehat \phi_h
&=
\frac{\sum_{m\in\mathcal{M}} Y_h(m)}
     {\sum_{m\in\mathcal{M}} \mathbf{1}\{\delta_m=1\}R_h(m)}.
\end{align*}
By construction,
\[
\widehat q_h=\widehat \pi_h\,\widehat \phi_h.
\]
Earlier drafts denoted \(\widehat q_h\) by \(\widehat p_h\). In the approved
architecture, \(\widehat q_h\) is the preferred notation and
\(\widehat p_{\mathrm{obs}}\) is reserved for the descriptive ever-reply share
\[
\widehat p_{\mathrm{obs}}
=
\frac{1}{|\mathcal{M}|}\sum_{m\in\mathcal{M}} \mathbf{1}\{\delta_m=1\}.
\]

All numerical readouts in this appendix come from the canonical
SimulaMet-backed evidence. A larger-archive Moltbook replication remains part
of the study design, but it is not yet a reported empirical result in the
current manuscript.

\subsection*{S2.1 Follow-up-standardized incidence by group}

Table \ref{tab:supp-horizon-incidence-groups} summarizes the primary
follow-up-standardized readouts at \(5\) minutes and \(1\) hour. Claimed versus
unclaimed is retained as descriptive heterogeneity only.

\begin{table}[h]
\centering
\caption{Follow-up-standardized incidence by group. The primary control-panel
throughput quantities are \(q_{5\mathrm{m}}\) and \(q_{1\mathrm{h}}\); earlier
drafts denoted the same quantities by \(p_h\). The in-window ever-reply share
\(p_{\mathrm{obs}}\) is descriptive only. Claimed/unclaimed excludes parents
with missing author identifier (\(n=906\)).}
\label{tab:supp-horizon-incidence-groups}
\small
\begin{tabular}{@{}llrrrr@{}}
\toprule
\textbf{Family} & \textbf{Group} & \textbf{Parents} & \textbf{\(q_{5\mathrm{m}}\) \%} & \textbf{\(q_{1\mathrm{h}}\) \%} & \textbf{\(p_{\mathrm{obs}}\) \%} \\
\midrule
overall & Overall & 223,316 & 9.42 & 9.82 & 9.60 \\
claimed\_status & Claimed & 20,667 & 18.95 & 19.56 & 19.23 \\
claimed\_status & Unclaimed & 201,743 & 8.48 & 8.86 & 8.65 \\
submolt\_category & Social/Casual & 189,765 & 10.80 & 11.21 & 10.95 \\
submolt\_category & Other & 16,396 & 1.98 & 2.29 & 2.26 \\
submolt\_category & Builder/Technical & 8,396 & 1.42 & 1.74 & 1.72 \\
submolt\_category & Philosophy/Meta & 5,831 & 1.32 & 1.70 & 1.80 \\
submolt\_category & Spam/Low-Signal & 2,495 & 0.88 & 0.93 & 0.88 \\
submolt\_category & Creative & 433 & 1.62 & 1.71 & 1.62 \\
\bottomrule
\end{tabular}
\end{table}

\section*{S3. Coverage-gap discipline}
\SupplementMixedEvidenceNote

The canonical comment stream contains one 41.7-hour interruption from
2026-01-31 10:37:53Z to 2026-02-02 04:20:50Z. The approved design therefore
treats gap handling as part of the identification strategy rather than as a
late robustness add-on.

\subsection*{S3.1 Gap-disambiguation evidence}

Table \ref{tab:supp-gap-disambiguation} shows that the interruption is
comment-stream-specific in the archive: comments disappear over the interval,
while posts, snapshots, and word-frequency records continue to arrive.

\begin{table}[h]
\centering
\caption{Comment-gap disambiguation evidence across raw archive tables. The
comment interval gap is 2026-01-31 10:37:53Z to 2026-02-02 04:20:50Z.}
\label{tab:supp-gap-disambiguation}
\small
\begin{tabular}{@{}lrrr@{}}
\toprule
\textbf{Table} & \textbf{Records} & \textbf{Records in comment-gap interval} & \textbf{Max inter-event gap (h)} \\
\midrule
comments & 226,173 & 0 & 41.72 \\
posts & 119,677 & 38,166 & 11.20 \\
snapshots & 114 & 39 & 2.28 \\
word\_frequency & 15,346 & 5,039 & 3.00 \\
\bottomrule
\end{tabular}
\end{table}

\subsection*{S3.2 Window robustness}

Table \ref{tab:supp-gap-window-robustness} reports the two-part decomposition
across the full window, contiguous pre-gap and post-gap windows, and the two
prespecified gap-overlap exclusions.

\begin{table}[h]
\centering
\caption{Two-part decomposition robustness across contiguous windows and
gap-overlap exclusions.}
\label{tab:supp-gap-window-robustness}
\scriptsize
\begin{tabular}{@{}lrrrrrrr@{}}
\toprule
\textbf{Scenario} & \textbf{Parents} & \textbf{Replies} & \textbf{Incidence \%} & \textbf{\(t_{50}\) (s)} & \textbf{\(t_{90}\) (s)} & \textbf{\(\Pr(\mathrm{reply}\le 30\mathrm{s})\) \%} & \textbf{\(\Pr(\mathrm{reply}\le 5\mathrm{m})\) \%} \\
\midrule
Full window & 223,316 & 21,430 & 9.60 & 4.55 & 50.05 & 8.47 & 9.41 \\
Pre-gap contiguous & 2,856 & 30 & 1.05 & 513.48 & 1699.77 & 0.04 & 0.35 \\
Post-gap contiguous & 220,460 & 21,400 & 9.71 & 4.55 & 48.66 & 8.58 & 9.53 \\
Exclude gap overlap (6h) & 220,460 & 21,400 & 9.71 & 4.55 & 48.66 & 8.58 & 9.53 \\
Exclude gap overlap (24h) & 220,460 & 21,400 & 9.71 & 4.55 & 48.66 & 8.58 & 9.53 \\
\bottomrule
\end{tabular}
\end{table}

\subsection*{S3.3 Horizon-standardized risk sets}

Table \ref{tab:supp-gap-horizon-standardized} reports the explicit risk-set
denominators used for the horizon-standardized readouts. Short-horizon values
remain stable, while the one-hour rate rises modestly under standardization,
consistent with right-censoring near the observation boundary rather than with a
slow conditional-timing regime.

\begin{table}[h]
\centering
\caption{Horizon-standardized reply probabilities using explicit risk sets.}
\label{tab:supp-gap-horizon-standardized}
\small
\begin{tabular}{@{}llrr@{}}
\toprule
\textbf{Scenario} & \textbf{Horizon} & \textbf{Risk-set \(n\)} & \textbf{\(\Pr(\mathrm{reply}\le t)\) \%} \\
\midrule
Full window & 30s & 223,312 & 8.47 \\
Full window & 5m & 223,102 & 9.42 \\
Full window & 1h & 217,472 & 9.82 \\
Pre-gap contiguous & 30s & 2,854 & 0.04 \\
Pre-gap contiguous & 5m & 2,842 & 0.35 \\
Pre-gap contiguous & 1h & 2,063 & 1.36 \\
Post-gap contiguous & 30s & 220,456 & 8.58 \\
Post-gap contiguous & 5m & 220,244 & 9.54 \\
Post-gap contiguous & 1h & 214,610 & 9.94 \\
Exclude gap overlap (6h) & 30s & 220,456 & 8.58 \\
Exclude gap overlap (6h) & 5m & 220,246 & 9.54 \\
Exclude gap overlap (6h) & 1h & 214,616 & 9.94 \\
Exclude gap overlap (24h) & 30s & 220,456 & 8.58 \\
Exclude gap overlap (24h) & 5m & 220,246 & 9.54 \\
Exclude gap overlap (24h) & 1h & 214,616 & 9.94 \\
\bottomrule
\end{tabular}
\end{table}

\section*{S4. Secondary diagnostic layers}
\SupplementMixedEvidenceNote

\subsection*{S4.1 Timing-model misspecification}

Parametric timing fits are retained only as diagnostics. Table
\ref{tab:supp-timing-model-fit} shows that even when event probability is
approximately calibrated, the fitted conditional timing quantiles materially
misstate the early-time mass that matters for the flagship control panel.

\begin{table}[h]
\centering
\caption{Observed versus fitted timing-model diagnostics for Moltbook. Large
early-quantile residuals indicate misspecification of the seconds-to-minutes
shape under a single-family parametric timing fit.}
\label{tab:supp-timing-model-fit}
\small
\begin{tabular}{@{}lrrr@{}}
\toprule
\textbf{Statistic} & \textbf{Observed} & \textbf{Fitted} & \textbf{Residual (fit - obs)} \\
\midrule
Event probability (\%) & 9.60 & 9.91 & +0.31 \\
\(p_{10}\) (s) & 3.67 & 6.00 & +2.33 \\
\(p_{50}\) (s) & 4.55 & 39.94 & +35.39 \\
\(p_{90}\) (s) & 50.05 & 134.98 & +84.93 \\
\bottomrule
\end{tabular}
\end{table}

\subsection*{S4.2 Model-to-observable coherence}

Table \ref{tab:supp-model-observable-validation} records the in-sample
coherence check used in the current manuscript. Incidence is approximately
matched, but non-root branching is underpredicted while depth tails are
overpredicted, consistent with strong depth dependence that a coarse homogeneous
fit does not capture.

\begin{table}[h]
\centering
\caption{Model-to-observable validation: predicted versus observed incidence,
non-root branching, and depth tails.}
\label{tab:supp-model-observable-validation}
\begingroup
\setlength{\tabcolsep}{4pt}
\scriptsize
\resizebox{\linewidth}{!}{%
\begin{tabular}{@{}lrrrrrrrr@{}}
\toprule
\textbf{Group} & \textbf{Pred. inc. \%} & \textbf{Obs. inc. \%} & \textbf{Pred. branch} & \textbf{Obs. branch} & \textbf{Pred. \(\Pr(D\ge 3)\)} & \textbf{Obs. \(\Pr(D\ge 3)\)} & \textbf{Pred. \(\Pr(D\ge 5)\)} & \textbf{Obs. \(\Pr(D\ge 5)\)} \\
\midrule
Overall & 9.91 & 9.60 & 0.104 & 0.134 & 0.0109 & 0.0013 & 0.00012 & 0.00001 \\
Claimed & 20.49 & 19.23 & 0.229 & 0.274 & 0.0526 & 0.0030 & 0.00276 & 0.00000 \\
Unclaimed & 8.90 & 8.65 & 0.093 & 0.120 & 0.0087 & 0.0012 & 0.00008 & 0.00001 \\
Builder/Technical & 1.72 & 1.72 & 0.017 & 0.017 & 0.0003 & 0.0001 & 0.00000 & 0.00000 \\
Creative & 1.62 & 1.62 & 0.016 & 0.016 & 0.0003 & 0.0000 & 0.00000 & 0.00000 \\
Other & 2.27 & 2.26 & 0.023 & 0.025 & 0.0005 & 0.0001 & 0.00000 & 0.00000 \\
Philosophy/Meta & 1.81 & 1.80 & 0.018 & 0.018 & 0.0003 & 0.0002 & 0.00000 & 0.00000 \\
Social/Casual & 11.36 & 10.95 & 0.121 & 0.154 & 0.0145 & 0.0015 & 0.00021 & 0.00001 \\
Spam/Low-Signal & 0.88 & 0.88 & 0.009 & 0.009 & 0.0001 & 0.0000 & 0.00000 & 0.00000 \\
\bottomrule
\end{tabular}
}
\endgroup
\end{table}

\subsection*{S4.3 Within-thread dependence sensitivity}

Table \ref{tab:supp-dependence-robustness} reports the one-parent-per-thread
sensitivity used to bound within-thread dependence. Incidence levels move more
than conditional timing, which is consistent with the interpretation that the
main bottleneck is whether a reply happens at all rather than how slow replies
are conditional on occurring.

\begin{table}[h]
\centering
\caption{One-parent-per-thread robustness against within-thread dependence.}
\label{tab:supp-dependence-robustness}
\small
\begin{tabular}{@{}lrrrr@{}}
\toprule
\textbf{Metric} & \textbf{Primary} & \textbf{One-parent/thread} & \textbf{Abs.\ diff.} & \textbf{Rel.\ diff. \%} \\
\midrule
Reply incidence \(\Pr(\delta=1)\) & 0.09596 & 0.07213 & 0.02383 & -24.84 \\
Conditional \(t_{50}\) (s) & 4.55 & 4.51 & 0.04 & -0.93 \\
Conditional \(t_{90}\) (s) & 50.05 & 41.08 & 8.97 & -17.92 \\
Kernel half-life (diagnostic; min) & 0.68451 & 0.43601 & 0.24850 & -36.30 \\
\bottomrule
\end{tabular}
\end{table}

\subsection*{S4.4 Periodicity details}

The longest contiguous segment contains 220,461 events over 63.5175 hours.
Modulo-4-hour event-time testing on that segment gives resultant length
\(r=0.0308\), Rayleigh \(Z=209.57\), Monte Carlo
\(p=5\times 10^{-6}\), and mean phase 153.2 minutes. This rejects exact
uniformity statistically, but the concentration is extremely small. In the same
evidence base, the four-hour target in the binned-count power spectral density
diagnostics yields first-order autoregressive null-calibrated \(p\)-values
0.508, 0.501, and 0.556 for 5-, 15-, and 30-minute bins, respectively.

\begin{table}[h]
\centering
\caption{Detectability mechanics for the modulo-4-hour periodicity diagnostic.}
\label{tab:supp-periodicity-detectability}
\small
\begin{tabular}{@{}ll@{}}
\toprule
\textbf{Component} & \textbf{Value} \\
\midrule
Target period \(\tau\) & 4.0 hours \\
Segment event count & 220,461 \\
Segment duration & 63.5175 hours \\
Null Monte Carlo reps (Rayleigh) & 200,000 \\
Power simulation method & \texttt{noncentral\_chi\_square\_monte\_carlo} \\
\(\kappa\) grid & 0.0, 0.2, \ldots, 3.0 \\
Power reps per \(\kappa\) & 50,000 \\
Seed & 20260208 \\
Estimated null size at \(\kappa=0\) & 0.04964 \\
First tested \(\kappa\) with estimated power \(\ge 80\%\) & 0.2 \\
Estimated power at \(\kappa=0.2\) & 1.0 \\
\bottomrule
\end{tabular}
\end{table}

Because the \(\kappa\) grid is coarse, the resulting \(\kappa^\star=0.2\) is a
grid-crossing reference rather than a sharp detectability threshold. The
corresponding mean resultant-length mapping is \(\rho\approx 0.0995\), which is
still well above the observed \(r=0.0308\).

\subsection*{S4.5 Topic keyword sensitivity}

Deterministic keyword remapping is used only to verify the direction of the
secondary topic contrast, not to estimate semantic effects. Across the baseline
and expanded trigger dictionaries, with and without exclusions for \emph{Other}
or small categories, the Social/Casual stratum remains higher-incidence and
faster-tailed than Philosophy/Meta. Because this layer is secondary and purely
diagnostic, the detailed variant-by-variant outputs are maintained in the
reproducibility pipeline rather than repeated as a load-bearing flagship table.


\bibliographystyle{plainnat}
\bibliography{paper/references}

\end{document}
